% ============================================================
\chapter*{Preface}
\addcontentsline{toc}{chapter}{Preface}
\markboth{Preface}{Preface}
% ============================================================

\section*{Course Objectives}

\textbf{Quantitative finance} lies at the intersection of mathematics,
statistics, and financial economics. This Master-level course aims to provide
students with a rigorous training covering the theoretical foundations and
numerical methods essential for understanding and practising modern
financial engineering.

Contemporary financial markets are characterised by increasing complexity:
exotic derivatives, high-frequency algorithmic trading, quantitative portfolio
management, and more recently the emergence of machine learning in financial
modelling. This course equips students with the mathematical and computational
tools necessary to navigate this environment.

The approach adopted is resolutely \textbf{mathematical and computational}:
each theoretical concept is accompanied by a Python implementation that
allows verification of analytical results and develops numerical intuition.

\section*{Prerequisites}

This course assumes that students possess solid knowledge in the following areas:

\begin{itemize}
    \item \textbf{Mathematical analysis}: differential and integral calculus,
          series, function spaces, ordinary differential equations.
    \item \textbf{Linear algebra}: vector spaces, matrices, eigenvalues,
          matrix decompositions (Cholesky, SVD).
    \item \textbf{Probability theory}: probability spaces, random variables,
          classical distributions, conditional expectation, limit theorems
          (LLN, CLT).
    \item \textbf{Statistics}: parametric estimation, hypothesis testing,
          linear regression, maximum likelihood estimation.
    \item \textbf{Programming}: proficiency in Python (or an equivalent
          scientific language), basic object-oriented programming.
\end{itemize}

\section*{Course Organisation}

The course is organised into eleven chapters, each designed to be covered
in one to two weeks of lectures accompanied by tutorials and computer
laboratory sessions.

\begin{description}
    \item[Chapter 1 --- Financial Markets and Modelling.]
        Introduction to markets, financial assets, modelling hypotheses,
        no-arbitrage principle, discrete-time models (Cox--Ross--Rubinstein).

    \item[Chapter 2 --- The Black--Scholes Model.]
        Model derivation, the Black--Scholes PDE, closed-form formulae
        for European options, computation and interpretation of the Greeks.

    \item[Chapter 3 --- Option Pricing Methods.]
        Binomial and trinomial trees, Monte Carlo methods, finite differences,
        variance reduction techniques.

    \item[Chapter 4 --- Hedging and Risk Management.]
        Delta hedging, dynamic hedging, gamma hedging, hedging costs,
        practical limitations of the Black--Scholes model.

    \item[Chapter 5 --- Stochastic Calculus in Finance.]
        Brownian motion, It\^o integral, It\^o's formula, Girsanov's theorem,
        martingale representation theorem.

    \item[Chapter 6 --- Stochastic Volatility.]
        Volatility smile and surface, the Heston model, the SABR model,
        calibration, local vs stochastic volatility.

    \item[Chapter 7 --- Interest Rate Models.]
        Term structure, Vasicek, CIR, Hull--White models, HJM framework,
        bond pricing and caps/floors.

    \item[Chapter 8 --- Portfolio Optimisation.]
        Markowitz theory, efficient frontier, CAPM, Black--Litterman model,
        robust optimisation.

    \item[Chapter 9 --- Risk Measures.]
        Value-at-Risk (VaR), Expected Shortfall (CVaR), coherence properties,
        computation methods, backtesting, stress testing.

    \item[Chapter 10 --- Exotic Derivatives.]
        Barrier options, Asian options, lookback options, basket options,
        analytical and numerical pricing.

    \item[Chapter 11 --- Machine Learning in Finance.]
        Neural networks for pricing, reinforcement learning for trading,
        anomaly detection, overfitting risks.
\end{description}

\section*{Pedagogical Methodology}

Each chapter follows a consistent structure:

\begin{enumerate}
    \item \textbf{Theory}: rigorous presentation of mathematical concepts
          with complete or sketched proofs.
    \item \textbf{Financial interpretation}: discussion of the economic
          significance of results and their limitations.
    \item \textbf{Implementation}: Python code illustrating numerical methods
          using the \texttt{numpy}, \texttt{scipy}, and \texttt{matplotlib}
          libraries.
    \item \textbf{Exercises}: theoretical and numerical problems of increasing
          difficulty to deepen understanding.
\end{enumerate}

\section*{Notation Conventions}

Throughout this course, we use the following notation:

\begin{center}
\begin{tabular}{cl}
    \toprule
    \textbf{Notation} & \textbf{Meaning} \\
    \midrule
    $(\Omega, \mathcal{F}, \PP)$ & Probability space \\
    $\E[X]$ & Expectation of the random variable $X$ \\
    $\E^\mathbb{Q}[\cdot]$ & Expectation under the risk-neutral measure \\
    $\Var(X)$ & Variance of $X$ \\
    $\Cov(X,Y)$ & Covariance of $X$ and $Y$ \\
    $\R, \R^+$ & Real numbers, positive reals \\
    $\N$ & Natural numbers \\
    $W_t$ or $B_t$ & Standard Brownian motion \\
    $(\mathcal{F}_t)_{t \geq 0}$ & Natural filtration \\
    $\mathbb{Q}$ & Risk-neutral probability measure \\
    $r$ & Risk-free interest rate (continuous) \\
    $\sigma$ & Volatility of the underlying asset \\
    $S_t$ & Price of the underlying asset at time $t$ \\
    $C(S,t), P(S,t)$ & European call and put prices \\
    $K$ & Strike price \\
    $T$ & Option maturity \\
    $\Delta, \Gamma, \Theta, \mathcal{V}, \rho$ & Greek letters (sensitivities) \\
    \bottomrule
\end{tabular}
\end{center}

\section*{Software and Tools}

The numerical implementations in this course use the scientific Python ecosystem:

\begin{itemize}
    \item \texttt{numpy}: matrix computation, linear algebra, random number
          generation.
    \item \texttt{scipy}: optimisation (\texttt{scipy.optimize}), integration,
          special functions, normal distribution (\texttt{scipy.stats.norm}).
    \item \texttt{matplotlib}: graphical visualisation, volatility surface
          plots, efficient frontiers.
    \item \texttt{pandas}: manipulation and analysis of financial time series.
    \item \texttt{scikit-learn}: classical machine learning algorithms
          (Chapter~11).
    \item \texttt{tensorflow} / \texttt{pytorch}: deep neural networks
          (Chapter~11).
\end{itemize}

\section*{Cross-Cutting Themes}

Several themes run through the entire course and deserve to be highlighted
from this preface:

\begin{itemize}
    \item \textbf{No-arbitrage}: the fundamental principle underpinning all
          pricing theory in quantitative finance.
    \item \textbf{Market completeness}: whether every contingent claim can be
          replicated by a self-financing trading strategy.
    \item \textbf{Model risk}: the consequences of using simplified mathematical
          models for complex financial realities.
    \item \textbf{Calibration vs estimation}: the distinction between calibrating
          a model to market prices and estimating its parameters from historical
          data.
\end{itemize}

\section*{References}

This course draws upon the following reference texts:

\begin{enumerate}
    \item \textsc{Shreve}, S.E. --- \textit{Stochastic Calculus for Finance
          I \& II}, Springer, 2004.
    \item \textsc{Hull}, J.C. --- \textit{Options, Futures, and Other Derivatives},
          Pearson, 2017 (10th edition).
    \item \textsc{Bj\"ork}, T. --- \textit{Arbitrage Theory in Continuous Time},
          Oxford University Press, 2009 (3rd edition).
    \item \textsc{Glasserman}, P. --- \textit{Monte Carlo Methods in Financial
          Engineering}, Springer, 2003.
    \item \textsc{Gatheral}, J. --- \textit{The Volatility Surface: A Practitioner's
          Guide}, Wiley, 2006.
    \item \textsc{Brigo}, D. \& \textsc{Mercurio}, F. --- \textit{Interest Rate
          Models --- Theory and Practice}, Springer, 2006.
    \item \textsc{McNeil}, A.J., \textsc{Frey}, R. \& \textsc{Embrechts}, P.
          --- \textit{Quantitative Risk Management: Concepts, Techniques and Tools},
          Princeton University Press, 2015.
    \item \textsc{L\'opez de Prado}, M. --- \textit{Advances in Financial Machine
          Learning}, Wiley, 2018.
    \item \textsc{Cont}, R. \& \textsc{Tankov}, P. --- \textit{Financial Modelling
          with Jump Processes}, Chapman \& Hall/CRC, 2004.
\end{enumerate}

\section*{Acknowledgements}

This course has benefited from the contributions of numerous colleagues
in mathematical finance. We are particularly grateful to students from
previous cohorts whose questions and pertinent remarks have considerably
improved the presentation of this material.

\bigskip
\begin{flushright}
    \textit{The Author}\\
    \textit{March 2026}
\end{flushright}

\newpage
